\documentclass[11pt]{article}

% \providecommand{\answerDirectory}{solutions}

../../macros/macros.tex

\inputAnswer{problem_0_answer}

\doHwHeader{Homework \BLUE{11}}

%%%%% EDIT THE FILES NAMED problem_X_answer.tex.  DON'T EDIT THIS FILE

%%%%%%%%%%%%%%%%%%%%%%%%%%%%%%%%%%%%%%%%%%%%%%%%%%%%%%%%%%%% 
%%%%%%%%%%%%%%%%%%%%%%%%%%%%%%%%%%%%%%%%%%%%%%%%%%%%%%%%%%%% 
%%%%%%%%%%%%%%%%%%%%%%%%%%%%%%%%%%%%%%%%%%%%%%%%%%%%%%%%%%%% 
%%%%%%%%%%%%%%%%%%%%%%%%%%%%%%%%%%%%%%%%%%%%%%%%%%%%%%%%%%%% 

\begin{document}

\gradescopeInit


%%%%% EDIT THE FILES NAMED problem_X_answer.tex.  DON'T EDIT THIS FILE 

\paragraph{General instructions.}
% First read \BLUE{Lecture Note 5}, then
Answer the problems given in the rest of this document.
Edit this LaTeX template to add your answers, as described in the file named \texttt{00\_README.txt},
then compile it to produce \BLUE{\texttt{homework\_11.pdf}}\ifGradescope{} for submission to gradescope\fi.

\ifGradescope
For this homework, when you submit to gradescope you don't need to identify the pages
for each problem.  We're using gradescope's ``fixed-length'' grading mode.
To make it work, your answer for each part should fit in the allotted space,
so that each subsequent answer is on the expected page in the expected location.
Your PDF should have
\BLUE{8}
pages.  
\fi

\paragraph{Don't know?  Say so.}
If you don't know the answer to a problem or any part of the problem,
we encourage you to simply write \emph{``I don't know''} instead of giving an answer.
(If you write ``I don't know'' you can also add any questions you might have.)

\paragraph{Gentle reminder: homework is for learning.}
As noted in \LN{prelim}, engaging fully with the homeworks is the best way to learn the material.
\emph{The primary role of homeworks is to help you learn the material,
  and the primary role of feedback on homeworks is to help you improve your ideas.}

Please try to explain your ideas as clearly as you can.
The clearer your explanations are, the more useful the feedback on them can be.

\paragraph{Collaboration and citation policy.}
If you are doing this homework as part of a course, follow the course's collaboration and citation policy.
In any case, at the end of each answer, list the people you collaborated with,
and cite every source that your answer uses ideas from
(other than the lecture notes and the sources listed in the lecture notes).
If there are none, write ``none''.

%%%%% EDIT THE FILES NAMED problem_X_answer.tex.  DON'T EDIT THIS FILE 

\vfill
\noindent{\footnotesize\copyrightNotice}

%%% Local Variables:
%%% mode: latex
%%% TeX-master: "homework_11"
%%% End:


%%%%%%%%%%%%%%%%%%%%%%%%%%%%%%%%%%%%%%%%%%%%%%%%%%%%%%%%%%%% 
%%%%%%%%%%%%%%%%%%%%%%%%%%%%%%%%%%%%%%%%%%%%%%%%%%%%%%%%%%%% 
%%%%%%%%%%%%%%%%%%%%%%%%%%%%%%%%%%%%%%%%%%%%%%%%%%%%%%%%%%%% 
%%%%%%%%%%%%%%%%%%%%%%%%%%%%%%%%%%%%%%%%%%%%%%%%%%%%%%%%%%%% 

%%%%% EDIT THE FILES NAMED problem_X_answer.tex.  DON'T EDIT THIS FILE 

%%%%%%%%%%%%%%%% PROBLEM 1  %%%%%%%%%%%%%%%%%

\newpage 


\problem\emph{(Sections 8.1 and 8.2)} \GradescopeComment{This should be on page 2 (top) of your PDF.}

Several families are having a dinner party.  To be more social, they would like to arrange seating so that no two members of the same family are at the same table.  There are $p$ families and $q$ tables.
For each $i\in\{1,2,\ldots,p\}$, family $i$ has $a_i$ members, where $a_i\le q$.
For each $j\in\{1,2,\ldots,q\}$, table $j$ has $b_j$ seats, where $b_j\le p$.
The goal is to determine if there exists a seating arrangement
(assigning each family member to a seat)
such that no two guests are assigned to the same seat,
and no two members of the same family sit at the same table.
Here is the formal definition of the problem, as a decision problem:

\begin{mylist}\em
\item[]\prob{Family Seating} \dotfill
\item[\bf input:] positive integers $p$, $q$, $a_1,\dots,a_p$, and $b_1,\dots,b_q$, where $\max_i a_i \le q$ and $\max_j b_j\le p$.
\item[\bf output:] `\textsf{Yes}' if it is possible to assign each guest to a seat such that (1) no two guests are assigned the same seat and (2) each table has at most one member of each family.
  If it is not possible, output `\textsf{no}'.
\end{mylist}

Recall the \prob{Max Flow} decision problem:
\begin{mylist}\em
% \item[]\prob{Max Flow} \dotfill
\item[\bf input:] digraph $G=(V,E)$ with edge capacities, source and sink $s, t\in V$, and a value $k\ge 0$
\item[\bf output:] `\textsf{Yes}' if $G$ has an $s$-$t$ flow of value at least $k$, otherwise `\textsf{no}'.
\end{mylist}

Develop a poly-time reduction $r$ from \prob{Family Seating} to \prob{Max Flow}.

\emph{Hint: First study the reduction from \prob{Bipartite Matching} to \prob{Max Flow}.}

\begin{enumerate}[label=(\alph*)]
\item 
  Define your reduction $r$
  by defining, for every possible instance $I=(p, q, a, b)$ of  \prob{Family Seating},
  what output $(G, s, t, k) = r(I)$ the reduction $r$ gives.
  Your definition must state precisely how $G$, $s$, $t$, and $k$ are defined for the given $(p, q, a, b)$.

\item
  Show that your reduction is a \emph{polynomial-time} reduction,
  by explaining carefully how, given $I$, one can compute the output $r(I)$ in time polynomial in the input size ($p+q$).

\item
  Consider the example input $I$ with $p=2$, $q=3$, $a=(2,2)$, $b=(1,2,2)$.
  Given this specific $I$, calculate the output $(G, s, t, k) = r(I)$.
  Draw the resulting flow network.  What is $k$? 

\item
For your reduction $r$ to be correct,
it must be that, \emph{for every  \prob{Family Seating} instance $I$,
the answer to $I$ is `\textsf{yes}' if and only if the answer to the \prob{Max Flow} decision problem $r(I)$ is `\textsf{yes}'.}

Show the ``if'' direction.
Given an arbitrary \prob{Family Seating} instance $I$,
let $G=(V,E)$, $s, t$, and $k$ be the output of your reduction $r$ given $I$.
Describe carefully how, given any $s$-$t$ flow $f$ in $G$ of value at least $k$,
one can obtain a valid seating arrangement for $I$.
(Assume that, if all edge capacities in $G$ are integers, then on each edge $(u,w)\in E$ the flow $f(u,w)$ is an integer.)

\item
For the input $I$ from Part (c), the answer is \textsf{yes},
so the flow network $(G, s, t)=r(I)$ should have an (integer) flow of value at least $k$.
Draw such an flow in your flow network.
What is the seating arrangement corresponding to that flow?

\item
Now show the ``only if'' direction.
Given an arbitrary \prob{Family Seating} instance $I$,
let $G=(V,E)$, $s, t$, and $k$ be the output of $r$ given $I$.
Describe carefully how, given any valid seating arrangement for $I$,
one can obtain an $s$-$t$ flow $f$ in $G$ of value at least $k$.
\end{enumerate}


\newpage

% \vfill

\setlength{\ITEMHT}{1.9in}
{% \injectEnumerate
\inputAnswer{problem_1_answer}
}

\newpage

%%%%%%%%%%%%%%%%%%%%%%%%%%%%%%%%%%%%%%%%%%%%%%%%%%%%%%%%%%%% 
%%%%%%%%%%%%%%%%%%%%%%%%%%%%%%%%%%%%%%%%%%%%%%%%%%%%%%%%%%%% 
%%%%%%%%%%%%%%%%%%%%%%%%%%%%%%%%%%%%%%%%%%%%%%%%%%%%%%%%%%%% 
%%%%%%%%%%%%%%%%%%%%%%%%%%%%%%%%%%%%%%%%%%%%%%%%%%%%%%%%%%%% 

%%%%% EDIT THE FILES NAMED problem_X_answer.tex.  DON'T EDIT THIS FILE 

%%%%%%%%%%%%%%%% PROBLEM 2 %%%%%%%%%%%%%%%%%

\problem\emph{(Section 8.2; Lecture Note 9)} \GradescopeComment{This should be on page 6 (top) of your PDF.}

\prob{$K$-Path} is the following problem:
\begin{mylist}\em
\item[\bf input:] A tuple $(G=(V,E), s, t, k)$,
  where $G$ is a directed graph with non-negative integer edge weights,
  $s,t\in V$ are two vertices, and $k$ is a non-negative integer.
\item[\bf output:]
  `Yes' if $G$ has an $s$-$t$ path $P$ with $\weight P$ equal to $k$, otherwise `no'.
\end{mylist}

Recall that \prob{Subset Sum} is the following problem:
\begin{mylist}\em
\item[\bf input:]
  Sequence $x=(x_1, x_2, \ldots, x_n)$ of non-negative integers and a non-negative integer $T$.
\item[\bf output:]
  `Yes' if there a subset $S\subseteq \{1,\ldots,n\}$ such that $\sum_{i\in S} x_i = T$, otherwise `no'.
\end{mylist}

Describe a polynomial-time reduction from \prob{Subset Sum} to \prob{$k$-Path}.

\begin{enumerate}[label=(\alph*)]

\item
  Define your reduction $r$
  by defining, for every possible instance $(x=(x_1, x_2, \ldots, x_n), T)$ of  \prob{Subset Sum},
  what output $r(x, T)$ the reduction gives.
  The output will be a tuple $(G=(V, E), s, t, k)$.
  Your definition must state precisely how $G$, $s$, $t$, and $k$ are defined for the given $(x, T)$.
  \emph{Hint: Construct $G$ so its $s$-$t$ paths correspond to the subsets $S\subseteq\{1,\ldots,n\}$.}

\item
  Show that your reduction is a \emph{polynomial-time} reduction,
  by explaining carefully how, given $(x, T)$, one can compute the output $r(x, T)=(G, s, t, k)$
  in time polynomial in the input size $\Theta(n\log T)$.

\item
  Consider the example input $(x, T)$ with $x=(5, 9, 10, 20)$ and $T=29$.
  Calculate the output $(G, s, t, k) = r(x, T)$ for this input.
  What is $k$? 
  Draw the resulting edge-weighted digraph.  

\item
  For your reduction $r$ to be correct,
  it must be that, \emph{for every  \prob{Subset Sum} instance $(x, T)$,
    the answer to $(x, T)$ is `\textsf{yes}' if and only if the answer to the \prob{$k$-Path} instance $r(x, T)$ is `\textsf{yes}'.}

  Show the ``if'' direction.
  Given an arbitrary \prob{Subset Sum} instance $(x, T)$,
  let $G$, $s$, $t$, and $k$ be the output of your reduction $r$ given $(x, T)$.  That is, $(G, s, t, k) = r(x, T)$.
  Describe carefully how, given any $s$-$t$ path $P$ in $G$ with $\weight P = k$,
  one could obtain a subset $S\subseteq \{1,\ldots,n\}$ such that $\sum_{i\in S} x_i = T$.

\item
  For the input $(x, T)$ from Part (c), the answer is `\textsf{yes}',
  so, for $(G, s, t, k)=r(x, T)$, there should be an $s$-$t$ path in $G$ of weight $k$.
  What is such a path $P$ in $G$?
  What is the subset $S\subseteq \{1,2,3,4\}$ corresponding to this path $P$?

\item
  Now show the ``only if'' direction.
  Given an arbitrary \prob{Subset Sum} instance $(x, T)$,
  let $G=(V, E)$, $s$, $t$, and $k$ be the output of your reduction $r$ given $(x, T)$.  That is, $(G, s, t, k) = r(x, T)$.
  Describe carefully how, given any subset $S\subseteq \{1,\ldots,n\}$ such that $\sum_{i\in S} x_i = T$,
  one could obtain from $S$ an $s$-$t$ path $P$ in $G$ with $\weight P = k$.

\end{enumerate}

% \setlength{\ITEMHT}{0.6in}
% {\injectEnumerate
\inputAnswer{problem_2_answer}
% }


%%%%%%%%%%%%%%%%%%%%%%%%%%%%%%%%%%%%%%%%%%%%%%%%%%%%%%%%%%%%   
%%%%%%%%%%%%%%%%%%%%%%%%%%%%%%%%%%%%%%%%%%%%%%%%%%%%%%%%%%%%   
%%%%%%%%%%%%%%%%%%%%%%%%%%%%%%%%%%%%%%%%%%%%%%%%%%%%%%%%%%%%   
%%%%%%%%%%%%%%%%%%%%%%%%%%%%%%%%%%%%%%%%%%%%%%%%%%%%%%%%%%%%   

%%%%%   EDIT THE FILES NAMED problem_X_answer.tex.  DON'T EDIT THIS FILE 

%%%%%%%%%%%%%%%%   PROBLEM 3 %%%%%%%%%%%%%%%%%

% \newpage 

% \problem \emph{(LN TODO Section TODO)} \GradescopeComment{This should be on page TODO (TODO) of your PDF.}

% \vfill
% \setlength{\ITEMHT}{2in}
% {\injectEnumerate
% \inputAnswer{problem_3_answer}
% }

\end{document}

%%% Local Variables:
%%% mode: latex
%%% TeX-master: t
%%% End:
